\section{Khi dự đoán chưa đủ}
\label{sec:ch01-khi-du-doan-chua-du}

\index{khả năng diễn giải|(}%
Một hệ thống học máy thường được đánh giá bằng đầu ra: dự đoán có đúng trên
dữ liệu kiểm tra không, lỗi có nằm trong mức chấp nhận được không. Cách đánh
giá này trả lời một câu hỏi hẹp về hiệu năng, nhưng không tự trả lời liệu hệ
thống có đang dùng một quy luật mà ta chấp nhận được, có còn phù hợp khi môi
trường thay đổi, hay có đẩy rủi ro sang một nhóm người mà dữ liệu kiểm tra
không làm lộ ra.

Nhu cầu về \tn{khả năng diễn giải}{interpretability} nảy sinh từ chính khoảng
trống đó. Tuy vậy, Lipton chỉ ra rằng từ này gộp nhiều mục đích khác nhau:
tạo niềm tin, tìm quan hệ nhân quả, chuyển giao sang tình huống mới, lấy thêm
thông tin cho người ra quyết định, và bảo đảm tính công bằng~\autocite{p13mythos}.
Gộp các mục đích ấy dưới một nhãn khiến một phương
pháp có thể được khen là \enquote{dễ hiểu} mà không ai nói rõ ai sẽ dùng nó,
để làm việc gì, và bằng tiêu chuẩn nào.

Xét một mô hình dự báo nguy cơ cho bệnh nhân. Một điểm số chính xác trên tập
kiểm tra chưa cho biết mô hình đang dùng tín hiệu bệnh lý hay đang tận dụng
một dấu vết của quy trình điều trị cũ. Bác sĩ có thể cần lời giải thích để
kiểm tra khả năng thứ hai, nhà nghiên cứu có thể cần nó để nêu một giả thuyết
về bệnh, còn người quản lý có thể cần nó để xem một tiêu chí công bằng đã bị
bỏ sót hay chưa. Như vậy, cùng một lời giải thích nhưng ba người dùng cần nó
cho ba việc kiểm tra khác nhau.

\index{hộp đen}%
Vì vậy, ta không bắt đầu bằng câu hỏi mô hình nào là \tn{hộp đen}{black box},
mà bằng câu
hỏi: điều gì còn nằm ngoài hàm mục tiêu và phép đo hiện có? Doshi-Velez và
Kim gọi tình trạng này là sự không đầy đủ của đặc tả bài toán: điều ta quan
tâm không thể liệt kê trọn vẹn để tối ưu hoặc để kiểm thử trực tiếp~\autocite{p14rigorous}.
Lời giải thích không tự lấp khoảng trống đó; nó chỉ đưa thêm bằng chứng để
con người kiểm tra khoảng trống ấy.

\index{khả năng diễn giải|)}%
